\documentclass[a4paper,12pt]{article}
\usepackage[utf8]{inputenc}
\usepackage[T1]{fontenc}
\usepackage[english]{babel}

% Questo comando importa i miei comandi personali che semplificano due o tre cagate
\usepackage{float}
\usepackage{url}
\usepackage{xcolor}
\usepackage{pdfpages}
\usepackage{graphicx}
\usepackage{mathtools,amssymb}
\usepackage{amsmath}

%Pro and cons
\usepackage{amssymb}% http://ctan.org/pkg/amssymb
\usepackage{pifont}% http://ctan.org/pkg/pifont
\newcommand{\cmark}{\ding{51}}%
\newcommand{\xmark}{\ding{55}}%
\newcommand{\pro}{\item[\cmark]}
\newcommand{\con}{\item[\xmark]}

\renewcommand{\tilde}{$\sim$}

%sets
\newcommand{\set}[1]{\{#1\}}

% Per numerare i paragrafi
\setcounter{tocdepth}{4}
\setcounter{secnumdepth}{4}
\let\oldpar\paragraph
\renewcommand{\paragraph}[1]{\oldpar{#1}\mbox{}\\}

%Per metter una freccia in modo fast
\newcommand{\fdx}{$\rightarrow$\ }

%Comando per creare nuove definizioni stile blocco
\newcommand{\definition}[2]{
	\begin{table}[H]
	\centering
		\begin{tabular}{|p{0.9\linewidth}}
		\textbf{#1}\\ %Titolo della definzione
		#2\\%Testo della definizione
		\end{tabular}
	\end{table}
	\noindent
}

%Comando esempio
\newcommand{\esempio}[1]{
	\begin{table}[H]
	\centering
		\begin{tabular}{p{0.9\linewidth}}
		\textbf{EX}\\ %Titolo della definzione
		\hline
		#1\\%Testo della definizione
		\end{tabular}
	\end{table}
	\noindent
}

\newcommand{\df}[2]{\definition{#1}{#2}}

\input{insbox.tex}
\newcommand{\lessonDate}[1]{\InsertBoxR{0}{\tiny{#1}}}

\newcommand{\E}{\`E\space}

\usepackage{listings}
\lstset{language=C++,
                keywordstyle=\color{blue},
                stringstyle=\color{red},
                commentstyle=\color{green},
                morecomment=[l][\color{magenta}]{\#}
}
\usepackage{hyperref}
\usepackage{cleveref}

% Tema scuro
%\usepackage{xcolor}
%
%\pagecolor[rgb]{0.25,0.25,0.25}
%\color[rgb]{1,1,1}
%\usepackage[demo]{graphicx}
\usepackage{caption}
\usepackage{subcaption}
\renewcommand\UrlFont{\color{blue}\rmfamily}
\usepackage{tikz}
\newcommand{\lira}{%
\begin{tikzpicture}[x=0.08em, y=0.08em, xscale=0.03, yscale=-0.03, inner sep=0pt, outer sep=0pt]
\fill (54.3355,9.3092) .. controls (70.3869,9.3092) and (79.7110,9.3092) ..
  (82.3075,9.3092) .. controls (82.3075,12.1418) and (82.3075,24.2985) ..
  (82.3075,45.7791) .. controls (82.3075,67.2598) and (82.3075,79.2984) ..
  (82.3075,81.8950) -- (167.2860,51.0903) .. controls (167.2860,59.3521) and
  (167.2860,66.7877) .. (167.2860,73.3971) -- (82.3075,104.2018) .. controls
  (82.3075,105.1460) and (82.3075,107.1525) .. (82.3075,110.2211) .. controls
  (82.3075,113.2898) and (82.3075,115.2962) .. (82.3075,116.2404) --
  (128.3375,99.9529) -- (167.2860,85.4358) .. controls (167.2860,86.8521) and
  (167.2860,90.4518) .. (167.2860,96.2351) .. controls (167.2860,102.0184) and
  (167.2860,105.5001) .. (167.2860,106.6804) .. controls (167.2860,107.6246) and
  (167.0499,108.0967) .. (166.5778,108.0967) -- (84.7861,137.4851) .. controls
  (83.8419,137.9572) and (83.0157,138.4293) .. (82.3075,138.9014) .. controls
  (82.3075,153.3005) and (82.3075,173.1878) .. (82.3075,198.5633) .. controls
  (82.3075,223.9388) and (82.3075,243.8262) .. (82.3075,258.2253) .. controls
  (82.3075,258.2253) and (82.3075,258.3433) .. (82.3075,258.5794) .. controls
  (82.3075,258.8154) and (82.3075,258.9334) .. (82.3075,258.9334) .. controls
  (103.7882,256.3369) and (122.7903,248.1931) .. (139.3139,234.5021) .. controls
  (157.4899,219.6309) and (169.6465,201.1009) .. (175.7838,178.9121) .. controls
  (178.3804,168.7619) and (179.6787,158.6117) .. (179.6787,148.4614) .. controls
  (179.6787,148.4614) and (189.1207,148.4614) .. (208.0048,148.4614) .. controls
  (208.0048,171.3584) and (202.4576,192.9571) .. (191.3632,213.2575) .. controls
  (183.5735,228.8369) and (172.9512,242.2918) .. (159.4963,253.6223) .. controls
  (146.9856,264.7167) and (132.9405,273.2145) .. (117.3611,279.1158) .. controls
  (97.0607,286.9055) and (76.0522,289.6201) .. (54.3355,287.2596) .. controls
  (54.3355,271.9163) and (54.3355,248.9013) .. (54.3355,218.2146) .. controls
  (54.3355,187.5279) and (54.3355,164.3949) .. (54.3355,148.8155) .. controls
  (54.3355,148.8155) and (52.8011,149.4057) .. (49.7325,150.5859) --
  (1.2239,167.9357) .. controls (1.2239,160.8541) and (1.2239,153.4185) ..
  (1.2239,145.6288) -- (54.3355,126.5087) .. controls (54.3355,125.0924) and
  (54.3355,120.9615) .. (54.3355,114.1160) .. controls (51.7389,115.2962) and
  (48.2571,116.7126) .. (43.8902,118.3649) .. controls (39.5232,120.0173) and
  (36.7496,120.9615) .. (35.5694,121.1975) -- (7.5973,131.4658) .. controls
  (7.5973,131.4658) and (6.8301,131.7018) .. (5.2958,132.1739) .. controls
  (3.7615,132.6460) and (2.4042,133.0001) .. (1.2239,133.2361) .. controls
  (1.2239,121.9057) and (1.2239,114.4701) .. (1.2239,110.9293) --
  (54.3355,92.1632) .. controls (54.3355,81.7770) and (54.3355,67.9680) ..
  (54.3355,50.7362) .. controls (54.3355,33.5045) and (54.3355,19.6955) ..
  (54.3355,9.3092) -- cycle;

\end{tikzpicture}}


\sloppy
\begin{document}

\begin{titlepage}
\begin{center}
	\Large{\textbf{Documentation for Algorithms for Massive Datasets project: Turkish lira recognizer}}
\vfill
\normalsize{Caccaro Sebastiano}\\
\normalsize{Cavagnino Matteo}\\
\normalsize{A.A.2019/2020}
\end{center}
\end{titlepage}
\pagenumbering{Roman}

\vspace*{\fill}
\textit{We declare that this material, which We now submit for assessment, is entirely our own work and has not been taken from the work of others, save and to the extent that such work has been cited and acknowledged within the text of our work. We understand that plagiarism, collusion, and copying are grave and serious offences in the university and accept the penalties that would be imposed should we engage in plagiarism, collusion or copying. This assignment, or any part of it, has not been previously submitted by us or any other person for assessment on this or any other course of study.}
\vspace*{\fill}

\newpage

\tableofcontents

\clearpage
%magari c'è qualcosa di un po piu slick di una section per questo

\pagenumbering{arabic}

\newpage
\section{Introduction}

The goal of this project is to build and train a classifier for Turkish Liras based on the \textit{Turkish Lira Banknote Dataset} (\autoref{sec:dataset}).\\
The classifier is developed using Convolutional Neural Networks (from now on CNN). The CNN is coded using python3 and using the Tensorflow and Keras libraries.

\subsection{References and sources}
In order to develop the project some other sources other than the ones given during the course were consulted: all the links are cited in the \textit{Reference} section of this document.\\
The sources include also some links from the \textit{Tensorflow official documentation} from which most of the code in the project has been adapted.

\subsection{Links}
All of the code from the project is available at the following link:
\begin{center}
\url{https://github.com/sebacaccaro/Progetto_AFMD}
\end{center}
The code is available both in a python file and in a notebook which is meant to be run on Google Colab. \\
Note that in both cases a Kaggle API key is needed in order download the dataset and run the code.






\newpage
\section{Dataset}
\label{sec:dataset}
The dataset used for the project is the \textit{Turkish Lira Banknote Dataset}, available at the following link:
\begin{center}
\url{https://www.kaggle.com/baltacifatih/turkish-lira-banknote-dataset}
\end{center}

The dataset is composed by 6000 images of Turkish Lira Banknotes, organized as follows:
\begin{itemize}
\item 1000 pictures of 5\lira\ banknotes
\item 1000 pictures of 10\lira\ banknotes
\item 1000 pictures of 20\lira\ banknotes
\item 1000 pictures of 50\lira\ banknotes
\item 1000 pictures of 100\lira\ banknotes
\item 1000 pictures of 200\lira\ banknotes
\end{itemize}

Images are already splitted in \textit{train} and \textit{validation}. The training set contains 925 images for each banknote, leading to 92.5\% of the images (5550) begin used for training and the remaining 7.5\% (450) being used for validation.

\subsection{Images}
All the pictures in the dataset consist of picture of banknotes taken with different perspectives and conditions: some are just laying on different surfaces, some are held in hand, blurred, folded or partially out of frame, ecc.

\begin{figure}[H]
\centering
\begin{subfigure}{.5\textwidth}
  \centering
  \includegraphics[width=1\linewidth]{Immagini/folded}
  \label{fig:sub1}
\end{subfigure}%
\begin{subfigure}{.5\textwidth}
  \centering
  \includegraphics[width=\textwidth]{Immagini/plain}
  \label{fig:sub2}
\end{subfigure}
\caption{Example of images of 20\lira\ banknotes }
\label{fig:test}
\end{figure}

The author of dataset already performed brightness adjustment, noise reduction and image flipping when necessary.\\
All images are in PNG format with a 1280 x 720 resolution.

\subsection{Preprocessing}
Taking into consideration what has already been said about the images, it is clear that there is no need to apply other forms of corrections. Therefore, the only steps necessary prior to the train phase are:
\begin{itemize}
	\item Image scaling and decoding
	\item Dataset loading and batch splitting
\end{itemize}

\subsubsection{Image Scaling}
All the images come in a 1280x70 resolution. Working with such high resolution it's not recommended when training a CNN as it would take a big tall on the training speed due to the increasing number of parameters and RAM usage. For this reason all image are reduced by scale factor of 5, bringing their resolution to 140x256.\\
Every image is also decoded in a 3D tensor with 3 channels corresponding to its RGB values.

\subsubsection{Dataset loading and batch splitting}
For both the training and the validation datasets a list is created containing the path of their images. Every element in the lists is then mapped to an image-label. The lists are then used to create the actual dataset variables.\\
Each dataset is divided in batches. After some tweaking, a number of 32 elements per batch proved to be  optimal for the considered application.\\
Each dataset is set to repeat, allowing for multiple epochs of training on the same batches. Moreover the two datasets are set to be prefetched: this allows the next batches to be loaded while the current one is being used, drastically reducing latency and improving throughput at the cost using additional memory.\\
The train dataset is also shuffled. This is not really necessary with the validation dataset, as it is only used for evaluation.


%esporre i modelli usati e i relativi risultati in ordine temporale in modo da fornire una sequenza dei ragionamenti fatti
\newpage
\section{Experiments and results}
Different models have been tested during the developing process; in this section some of those will be shown and the relative results will be discussed.
\subsection{Model summary}
For each tested model the following data will be reported:
\begin{itemize}
\item The NN architecture
\item Hyperparameters used
\item Data on accuracy for three repeated runs
\item Graph of one of the runs
\item Comment on the architecture and results
\end{itemize}
In the layer tables the input  layer will not be reported, as it always corresponds to resized image size (\texttt{144,256,3}).\\
Also note that some abbreviations are used in the Layer Config field in order to for the table to fit:
\begin{itemize}
\item \texttt{k} stands for \texttt{kernel size}
\item \texttt{s} stands for \texttt{strides}
\item \texttt{f} stands for \texttt{filters}
\item \texttt{p} stands for \texttt{pool size}
\item \texttt{r} stands for \texttt{rate}
\end{itemize}

%Commands needed for table
\newcommand{\conv}{Convolution(\texttt{Conv2d})}
\newcommand{\convP}[3]{\texttt{k=#1, s=#2, f=#3}}
\newcommand{\convKSF}[3]{\convP{#1}{#2}{#3}}

\newcommand{\flt}{Flatten(\texttt{Flatten})}

\newcommand{\dns}{Dense(\texttt{Dense})}
\newcommand{\dnsP}[1]{\texttt{u=#1}}

\newcommand{\pool}{MaxPooling(\texttt{MaxPooling2D})}
\newcommand{\poolN}{\texttt{p=2x2}}

\newcommand{\drop}{Dropout(\texttt{Dropout})}
\newcommand{\dropR}[1]{\texttt{r=#1}}

\newcommand{\bat}{Batch Norm.(\texttt{BatchN.})}
\subsection{Models}

\subsubsection{Baseline Model}
\begin{table}[H]
    \centering
	\begin{tabular}{lcccc}
	\textbf{Layer Type} & \textbf{Layer Config} & \textbf{Activation}  & \textbf{Output} & \textbf{Params}\\ \hline
	\conv	& \convP{5}{3}{5}	& relu		& \texttt{48,86,5} 	& \texttt{380}\\
	\flt		& /					& relu		& \texttt{20640}		& \texttt{0}\\
	\dns		& \dnsP{64}			& relu		& \texttt{64}		& \texttt{1321024}\\
	\dns		& \dnsP{6}			& softmax	& \texttt{64}		& \texttt{390}\\
	\end{tabular}
	%Total params: 1,321,794
	%Trainable params: 1,321,794
	%Non-trainable params: 0
\end{table}

\begin{table}[H]
	\centering
	\begin{tabular}{lc}
	\textbf{Param} & \textbf{Value}\\ \hline
	Batch Size 	& 32 \\
	Optimizer 	& Adam \\
	Base lr		& 0.001 \\
	\end{tabular}
\end{table}


\begin{figure}[H]
	\begin{center}
	\includegraphics[width=\linewidth]{Immagini/Baseline-1}
	\caption{Graph of the first run}
	\end{center}
\end{figure}
\begin{table}[H]
	\centering
	\begin{tabular}{cccccc}
		\textbf{Run} &\textbf{Loss}&\textbf{V.Loss} &\textbf{Acc.}&\textbf{V.Acc.}&\textbf{$\Delta$ Acc.} \\ \hline
		1	& 0.0016		& 0.0530		& 1.0000		& 0.9754		& 0.0246 \\
		2	& 0.0012		& 0.0342		& 1.0000		& 0.9821		& 0.0179 \\
		3	& 0.0060		& 0.0497		& 0.9996		& 0.9866		& 0.0130 \\
		\textbf{Avg} & \textbf{0.0029}	& \textbf{0.0456}	& \textbf{0.9996} 	& \textbf{0.9814}	& \textbf{0.0185} 
	\end{tabular}
\end{table}




\subsubsection{Convolution Model}
\begin{table}[H]
    \centering
	\begin{tabular}{lcccc}
	\textbf{Layer Type} & \textbf{Layer Config} & \textbf{Activation}  & \textbf{Output} & \textbf{Params}\\ \hline
	\conv	& \convKSF{5}{3}{5}	& relu		& \texttt{48,86,5} 	& \texttt{380}\\
	\conv	& \convKSF{5}{2}{8}	& relu		& \texttt{24,43,8} 	& \texttt{1008}\\	
	\conv	& \convKSF{3}{1}{12}	& relu		& \texttt{24,43,12} 	& \texttt{876}\\
	\conv	& \convKSF{3}{1}{15}	& relu		& \texttt{24,43,15} 	& \texttt{1635}\\
	\conv	& \convKSF{3}{1}{18}	& relu		& \texttt{24,43,18} 	& \texttt{2448}\\
	
	\flt		& /					& /		& \texttt{20640}		& \texttt{0}\\
	\dns		& \dnsP{64}			& relu		& \texttt{64}		& \texttt{1188928}\\
	\dns		& \dnsP{6}			& softmax	& \texttt{6}		& \texttt{390}\\
	\multicolumn{4}{r}{\textbf{TOTAL}}&{\textbf{1,195,665}}\\
	\end{tabular}
	%Total params: 1,195,665
	%Trainable params: 1,195,665
	%Non-trainable params: 0
\end{table}


\begin{table}[H]
	\centering
	\begin{tabular}{lc}
	\textbf{Param} & \textbf{Value}\\ \hline
	Batch Size 	& 32 \\
	Optimizer 	& Adam \\
	Base lr		& 0.001 \\
	Epochs		& 20 \\
	\end{tabular}
\end{table}


\begin{figure}[H]
	\begin{center}
	\includegraphics[width=\linewidth]{Immagini/conv-3}
	\caption{Graph of the third run}
	\end{center}
\end{figure}
\begin{table}[H]
	\centering
	\begin{tabular}{cccccc}
		\textbf{Run} &\textbf{Loss}&\textbf{V.Loss} &\textbf{Acc.}&\textbf{V.Acc.}&\textbf{$\Delta$ Acc.} \\ \hline
		1   & 1.0908e-04    &   0.0439  & 1.0000    & 0.9866    & 0.0134 \\
		2   & 8.7761e-05    &   0.0912  & 1.0000    & 0.9888    & 0.0112 \\
		3   & 1.8083e-04    &   0.0609  & 1.0000    & 0.9911    & 0.0089 \\
		\textbf{Avg} & \textbf{1.2589e-04} & \textbf{0.0653}	& \textbf{1.0000}	& \textbf{0.9889} 	& \textbf{0.0112} 
	\end{tabular}
\end{table}

This models is an iteration of the baseline model, achieved by adding four more convolution. The convolution have been added with the following logic in mind:
\begin{itemize}
\item The first two convolutions have a stride greater than 1, and that leads to a minimal reductions in the size of the tensor. This in turn should limit the number of parameters needed for training and in turn reduce overfitting.
\item The layers are set in way to progressively increase the number of channel and reduce height and width.
\end{itemize}
The results show the model is slightly less overfitting than the baseline. Both models reached 1.0000 accuracy on the training set, so this translate in an increased validation accuracy.
%NOTARE CHE AGGIUNGERE PIU LIVELLI NON AIUTA





\subsubsection{Convolution and Pooling Model}
\begin{table}[H]
    \centering
	\begin{tabular}{lcccc}
	\textbf{Layer Type} & \textbf{Layer Config} & \textbf{Activation}  & \textbf{Output} & \textbf{Params}\\ \hline
	\conv	& \convKSF{5}{1}{5}	& relu		& \texttt{144,256,5} 	& \texttt{380}\\
	\pool	& \poolN				&	/		& \texttt{72,128,8}		& 0	\\
	\conv	& \convKSF{5}{1}{8}	& relu		& \texttt{72,128,8} 		& \texttt{1008}\\
	\pool	& \poolN				&	/		& \texttt{36,64,12}		& 0	\\	
	\conv	& \convKSF{3}{1}{12}	& relu		& \texttt{36,64,12} 		& \texttt{876}\\
	\pool	& \poolN				&	/		& \texttt{18,32,15} 		& 0	\\
	\conv	& \convKSF{3}{1}{15}	& relu		& \texttt{18,32,15} 		& \texttt{1635}\\
	\pool	& \poolN				&	/		& \texttt{9,16,18}		& 0	\\
	\conv	& \convKSF{3}{1}{18}	& relu		& \texttt{9,16,18} 		& \texttt{2448}\\
	
	\flt		& /					& /			& \texttt{2592}			& \texttt{0}\\
	\dns		& \dnsP{64}			& relu		& \texttt{64}			& \texttt{165952}\\
	\dns		& \dnsP{6}			& softmax	& \texttt{6}				& \texttt{390}\\
	\multicolumn{4}{r}{\textbf{TOTAL}}&{\textbf{172,689}}\\
	\end{tabular}
	%Total params: 172,689
	%Trainable params: 172,689
	%Non-trainable params: 0
\end{table}


\begin{table}[H]
	\centering
	\begin{tabular}{lc}
	\textbf{Param} & \textbf{Value}\\ \hline
	Batch Size 	& 32 \\
	Optimizer 	& Adam \\
	Base lr		& 0.001 \\
	Epochs		& 20 \\
	\end{tabular}
\end{table}


\begin{figure}[H]
	\begin{center}
	\includegraphics[width=\linewidth]{Immagini/conv-pool-1}
	\caption{Graph of the first run}
	\end{center}
\end{figure}
\begin{table}[H]
	\centering
	\begin{tabular}{cccccc}
		\textbf{Run} &\textbf{Loss}&\textbf{V.Loss} &\textbf{Acc.}&\textbf{V.Acc.}&\textbf{$\Delta$ Acc.} \\ \hline
		1   & 1.8558e-04 &  0.0429  & 1.0000    & 0.9933    & 0.0067\\
		2   & 6.4439e-04 &  0.0461  & 1.0000    & 0.9866    & 0.0134\\
		3   & 1.2224e-04 &  0.0133  & 1.0000    & 0.9933    & 0.0067\\
		\textbf{Avg} & \textbf{3.1740e-04} &  \textbf{0.0341}  & \textbf{1.0000}    & \textbf{0.9911}    & \textbf{0.0893}
	\end{tabular}
\end{table}

This model applies the use of the Conv-Pool technique in order to convolve and reduce the data at every step. Since the reduction is performed by various MaxPooling layers, the first two convolutions have now a stride equal to 1. Pooling helps by reducing the tensor dimension while retaining most of the information. 
Moreover, the model has a number of parameters which are an order of magnitude lower than the previous model. Therefore the model has now little to no overfitting and has gained a good amount of validation accuracy.

\subsubsection{Convolution and Pooling Model with Dropout}
%Questo modello fa più schifo di quelli di prima, dropout scartato
\begin{table}[H]
    \centering
	\begin{tabular}{lcccc}
	\textbf{Layer Type} & \textbf{Layer Config} & \textbf{Activation}  & \textbf{Output} & \textbf{Params}\\ \hline
	\conv	& \convKSF{5}{1}{5}	& relu		& \texttt{144,256,5} 	& \texttt{380}\\
	\pool	& \poolN				&	/		& \texttt{72,128,8}		& 0	\\
	\conv	& \convKSF{5}{1}{8}	& relu		& \texttt{72,128,8} 		& \texttt{1008}\\
	\pool	& \poolN				&	/		& \texttt{36,64,12}		& 0	\\	
	\conv	& \convKSF{3}{1}{12}	& relu		& \texttt{36,64,12} 		& \texttt{876}\\
	\pool	& \poolN				&	/		& \texttt{18,32,15} 		& 0	\\
	\conv	& \convKSF{3}{1}{15}	& relu		& \texttt{18,32,15} 		& \texttt{1635}\\
	\pool	& \poolN				&	/		& \texttt{9,16,18}		& 0	\\
	\conv	& \convKSF{3}{1}{18}	& relu		& \texttt{9,16,18} 		& \texttt{2448}\\
	\drop	& \dropR{0.75}		&	/		& \texttt{9,16,18}		& 0\\
	
	\flt		& /					& /			& \texttt{2592}			& \texttt{0}\\
	\dns		& \dnsP{64}			& relu		& \texttt{64}			& \texttt{165952}\\
	\drop	& \dropR{0.75}		&	/		& \texttt{64}			& 0\\
	\dns		& \dnsP{6}			& softmax	& \texttt{6}				& \texttt{390}\\
	\end{tabular}
	%Total params: 172,689
	%Trainable params: 172,689
	%Non-trainable params: 0
\end{table}


\begin{table}[H]
	\centering
	\begin{tabular}{lc}
	\textbf{Param} & \textbf{Value}\\ \hline
	Batch Size 	& 32 \\
	Optimizer 	& Adam \\
	Base lr		& 0.001 \\
	Epochs		& 20 \\
	\end{tabular}
\end{table}


\begin{figure}[H]
	\begin{center}
	\includegraphics[width=\linewidth]{Immagini/conv-drop-pool-1}
	\caption{Graph of the first run}
	\end{center}
\end{figure}
\begin{table}[H]
	\centering
	\begin{tabular}{cccccc}
		\textbf{Run} &\textbf{Loss}&\textbf{V.Loss} &\textbf{Acc.}&\textbf{V.Acc.}&\textbf{$\Delta$ Acc.} \\ \hline
		1   & 0.0582    & 0.0666    & 0.9817    & 0.9821    & -0.0004\\
		2   & 0.0063    & 0.0269    & 0.9987    & 0.9933    & 0.0054\\
		3   & 0.1172    & 0.0615    & 0.9654    & 0.9866    & -0.0212\\
		\textbf{Avg} & \textbf{0.0606}    & \textbf{0.0517}    & \textbf{0.9820}    & \textbf{0.9873}    & \textbf{-0.0054}
	\end{tabular}
\end{table}


\subsubsection{Batch Normalization Model}
\begin{table}[H]
    \centering
	\begin{tabular}{lcccc}
	\textbf{Layer Type} & \textbf{Layer Config} & \textbf{Activation}  & \textbf{Output} & \textbf{Params}\\ \hline
	\conv	& \convKSF{5}{1}{5}	& 	/		& \texttt{144,256,5} 	& \texttt{375}\\
	\pool	& \poolN				&	/		& \texttt{72,128,8}		& 0	\\
	\bat		& 		/			&	/		& \texttt{72,128,8}		& 15 \\
	\multicolumn{5}{c}{\texttt{Relu Activation}}\\
	
	\conv	& \convKSF{5}{1}{8}	& 	/		& \texttt{72,128,8} 		& \texttt{1000}\\
	\pool	& \poolN				&	/		& \texttt{36,64,12}		& 0	\\	
	\bat		& 		/			&	/		& \texttt{36,64,12}		& 24 \\
	\multicolumn{5}{c}{\texttt{Relu Activation}}\\	
	
	\conv	& \convKSF{3}{1}{12}	& 	/		& \texttt{36,64,12} 		& \texttt{864}\\
	\pool	& \poolN				&	/		& \texttt{18,32,15} 		& 0	\\
	\bat		& 		/			&	/		& \texttt{18,32,15} 		& 36 \\
	\multicolumn{5}{c}{\texttt{Relu Activation}}\\
	
	\conv	& \convKSF{3}{1}{15}	& 	/		& \texttt{18,32,15} 		& \texttt{1620}\\
	\pool	& \poolN				&	/		& \texttt{9,16,18}		& 0	\\
	\bat		& 		/			&	/		& \texttt{9,16,18} 		& 45 \\
	\multicolumn{5}{c}{\texttt{Relu Activation}}\\
	\drop	& \dropR{0.06}		&	/		& \texttt{9,16,18}		& 0\\
	
	\conv	& \convKSF{3}{1}{18}	& 	/		& \texttt{9,16,18} 		& \texttt{2430}\\	
	\flt		& /					& 	/		& \texttt{2592}			& \texttt{0}\\
	\bat		& 		/			&	/		& \texttt{2592}	 		& 7776 \\
	\multicolumn{5}{c}{\texttt{Relu Activation}}\\	
	
	\dns		& \dnsP{64}			& 	/		& \texttt{64}			& \texttt{165888}\\
	\bat		& 		/			&	/		& \texttt{64}	 		& 192 \\
	\multicolumn{5}{c}{\texttt{Relu Activation}}\\
	
	\drop	& \dropR{0.06}		&	/		& \texttt{64}		& 0\\
	
	\dns		& \dnsP{6}			& softmax	& \texttt{6}			& \texttt{390}\\
	\end{tabular}
%Total params: 180,655
%Trainable params: 175,263
%Non-trainable params: 5,392
\end{table}


\begin{table}[H]
	\centering
	\begin{tabular}{lc}
	\textbf{Param} & \textbf{Value}\\ \hline
	Batch Size 	& 32 \\
	Optimizer 	& Adam \\
	Base lr		& 0.00005 \\
	Epochs		& 50 \\
	\end{tabular}
\end{table}


\begin{figure}[H]
	\begin{center}
	\includegraphics[width=\linewidth]{Immagini/bn}
	\caption{Graph of the third run}
	\end{center}
\end{figure}
\begin{table}[H]
	\centering
	\begin{tabular}{cccccc}
		\textbf{Run} &\textbf{Loss}&\textbf{V.Loss} &\textbf{Acc.}&\textbf{V.Acc.}&\textbf{$\Delta$ Acc.} \\ \hline
	1   & 0.0276    & 0.0439    & 0.9962    & 0.9933    & 0.0029\\
	2   & 0.0255    & 0.0510    & 0.9962    & 0.9888    & 0.0074\\
	3   & 0.0227    & 0.0353    & 0.9969    & 0.9933    & 0.0036\\
	\textbf{Avg} & \textbf{0.0253}    & \textbf{0.0434}    & \textbf{0.9964}    & \textbf{0.9918}    & \textbf{0.0046}\\ 
	\end{tabular}
\end{table}






\subsection{Models comparison}

When deciding which is best, the following factors should be taken into account:
\begin{itemize}
\item \textbf{Accuracy}: the more accurate the model, the better. The measure to look for is validation accuracy.
\item \textbf{Generalization}: the model needs to work with images outside of the training dataset. If the model generalizes well, the difference between its accuracy and validation accuracy should be little to none.
\item \textbf{Number of parameters:} if two models have similar performances, the one with less parameters is to be preferred. Having less parameters generally makes the model faster in training and evaluation. It also comes with the added benefit of having a lower chance of overfitting.
\end{itemize}

\begin{table}[H]
	\centering
	\begin{tabular}{ccc cccc}
	\textbf{\#} & \textbf{Loss} & \textbf{V.Loss} & \textbf{Acc.} & \textbf{V.Acc} & \textbf{$\Delta$ Acc.} & \textbf{Params.}\\ \hline
	1 & 0.0029 & 0.0456 & 0.9996 & 0.9814 & 0.0185 & 1,321,794\\ 
	2 & 0.0001 & 0.0653 & 1.0000	& 0.9889 & 0.0112 & 1,195,665\\
	3 & 0.0003 & 0.0341 & 1.0000 & 0.9911 & 0.0893 & 172,689\\
	4 & 0.0606 & 0.0517 & 0.9820 & 0.9873 & -0.0054 & 172,689\\
	5 & 0.0253 & 0.0434 & 0.9964 & 0.9918 & 0.0046 & 180,655\\
	\end{tabular}
	\caption{Average performance of the various model}
\end{table}

On this grounds it's clear that models 4 and 5 are the best candidates for the given dataset.






\newpage
\section{Conclusions}
As seen in the previous sections, after a few experiments, this project has been succesful in classifying the images of Turkish 
Lira banknotes with high accuracy and low data loss.
(altro?) aggiungere migliori risultati ottenuti - 




Esempio di citazione\cite{latexcompanion}

\newpage
\addcontentsline{toc}{section}{References}
\begin{thebibliography}{9}
\bibitem{kaggle-colab} 
Easiest way to download kaggle data in Google Colab\\
\url{https://www.kaggle.com/general/74235}

\bibitem{images} 
Load images\\
\url{https://www.tensorflow.org/tutorials/load_data/images}

\bibitem{images} 
Classification\\
\url{https://www.tensorflow.org/tutorials/images/classification}

\bibitem{CNN}
Convolutional Neural Network (CNN)\\
\url{https://www.tensorflow.org/tutorials/images/cnn} 

\bibitem{CNNGuide}
A Comprehensive Guide to Convolutional Neural Networks — the ELI5 way\\
\url{https://towardsdatascience.com/a-comprehensive-guide-to-convolutional-\\neural-networks-the-eli5-way-3bd2b1164a53} 

\bibitem{CNNGuide2}
A guide to an efficient way to build neural network architectures- Part II: Hyper-parameter selection and tuning for Convolutional Neural Networks using Hyperas on Fashion-MNIST\\
\url{https://towardsdatascience.com/a-guide-to-an-efficient-way\\-to-build-neural-network-architectures-part-ii-hyper-parameter-42efca01e5d73} 

\bibitem{notebook}
Prof. Notebook on deep learning\\







\end{thebibliography}


\end{document}
